\begin{abstract}
Finite-aperture scattering data can rule out an entire dielectric inclusion class, although they may leave individual pixels undetermined. We derive two results for the discrete nonlinear transverse magnetic volume integral equation. First, a family of material intervals that changes on a fixed set of cells admits a common exact reduction. Retaining the current and total-field traces there recovers the inertia of every data test up to a known positive block, including singular cases. Second, when the old readings exhaust the available negative directions, every compatible material satisfies a conditional quadratic bound on its response to future illumination. A separate index computed from the old sources determines whether this bound is an ellipsoid. In 24 material pairs the full and reduced tests agreed on all 2688 material hypotheses. Imposing exact laws on selected unknown cells gave no additional exclusions beyond exact elimination of known vacuum. On a 9216-cell model the data ruled out 247 of 256 positions for a core of specified size, with one verified compatible position and eight unresolved. In six finite-aperture examples a single additional beam separated two previously compatible classes. A uniform angle also achieved this result. The conditional bound establishes when a future beam suffices before its response is measured. On the native array a physically compatible class leaves seven receiving directions uncontrolled at saturation even with 32 illuminations. For a narrower class, eleven coherent beams above a calculated energy threshold yield a bounded receiving prediction; eleven angular beams at the same energy leave five directions uncontrolled.
\end{abstract}
